% =====================================================================
\chapter{Identidades trigonométricas}
\label{tema:2}
% =====================================================================

Una identidad trigonométrica es una igualdad entre expresiones trigonométricas que se cumple para todos los valores de la variable en que ambos miembros están definidos. Este tema demuestra las identidades fundamentales a partir de la fórmula del coseno de una diferencia, que a su vez se obtiene de una distancia en la circunferencia unidad, y muestra cómo usarlas para simplificar y verificar expresiones. Se sigue el capítulo~9 de~\cite{osat} (§9.1 a §9.4), equivalente al capítulo~7 de~\cite{osprecalc}; en el \cref{tema:6} se verá una segunda demostración de las fórmulas de la suma, mediante la exponencial compleja, que es la que propone el curso 18.03 del MIT~\cite[§6.4]{mit1803}.

\section{Identidades pitagóricas y verificación de identidades}

\begin{proposicion}[Identidades pitagóricas]\label{prop:pitagoricas}
Para todo $\theta$ en que estén definidas las funciones,
\[
\sin^{2}\theta+\cos^{2}\theta=1,\qquad 1+\tan^{2}\theta=\sec^{2}\theta,\qquad 1+\cot^{2}\theta=\csc^{2}\theta.
\]
\end{proposicion}

Para \emph{verificar} una identidad se transforma uno de sus miembros, mediante identidades ya demostradas y operaciones algebraicas reversibles, hasta obtener el otro. No es válido manipular la igualdad que se quiere probar como si fuera cierta (por ejemplo, elevando ambos miembros al cuadrado), porque así se pueden «demostrar» igualdades falsas~\cite[§9.1]{osat}.

\addtocounter{definicion}{1}%

\section{Fórmulas de la suma y de la diferencia}

\begin{teorema}[Coseno de una diferencia]\label{teo:cos-diferencia}
Para todos $\alpha,\beta\in\R$,
\[
\cos(\alpha-\beta)=\cos\alpha\cos\beta+\sin\alpha\sin\beta .
\]
\end{teorema}

\begin{corolario}[Fórmulas de la suma y la diferencia]\label{cor:suma}
Para todos $\alpha,\beta\in\R$,
\begin{align*}
\cos(\alpha\pm\beta)&=\cos\alpha\cos\beta\mp\sin\alpha\sin\beta,\\
\sin(\alpha\pm\beta)&=\sin\alpha\cos\beta\pm\cos\alpha\sin\beta,\\
\tan(\alpha\pm\beta)&=\frac{\tan\alpha\pm\tan\beta}{1\mp\tan\alpha\tan\beta}\quad\text{(cuando todas las tangentes están definidas).}
\end{align*}
\end{corolario}

\addtocounter{definicion}{1}%

\begin{corolario}[Composición de giros]\label{cor:giros}
Sea $R(\alpha)=\begin{pmatrix}\cos\alpha&-\sin\alpha\\ \sin\alpha&\cos\alpha\end{pmatrix}$. Entonces $R(\alpha)R(\beta)=R(\alpha+\beta)$ para todos $\alpha,\beta$.
\end{corolario}

La matriz $R(\alpha)$ lleva el vector $(\cos\theta,\sin\theta)$ al vector $(\cos(\theta+\alpha),\sin(\theta+\alpha))$, como se ve aplicando el \cref{cor:suma}: es el giro de ángulo $\alpha$. El \cref{cor:giros} dice que girar $\beta$ y después $\alpha$ equivale a girar $\alpha+\beta$.

\section{Ángulo doble y ángulo mitad}

\begin{proposicion}[Ángulo doble]\label{prop:doble}
Para todo $\alpha$,
\[
\sin 2\alpha=2\sin\alpha\cos\alpha,\qquad \cos 2\alpha=\cos^{2}\alpha-\sin^{2}\alpha=2\cos^{2}\alpha-1=1-2\sin^{2}\alpha,
\]
y $\tan 2\alpha=\dfrac{2\tan\alpha}{1-\tan^{2}\alpha}$ cuando los dos miembros están definidos.
\end{proposicion}

\begin{corolario}[Reducción de potencias y ángulo mitad]\label{cor:mitad}
Para todo $\alpha$,
\[
\cos^{2}\alpha=\frac{1+\cos 2\alpha}{2},\qquad \sin^{2}\alpha=\frac{1-\cos 2\alpha}{2}.
\]
En consecuencia, $\cos\frac{\alpha}{2}=\pm\sqrt{\frac{1+\cos\alpha}{2}}$ y $\sin\frac{\alpha}{2}=\pm\sqrt{\frac{1-\cos\alpha}{2}}$, donde el signo es el de la función en el cuadrante de $\alpha/2$, y $\tan\frac{\alpha}{2}=\frac{1-\cos\alpha}{\sin\alpha}=\frac{\sin\alpha}{1+\cos\alpha}$ cuando los denominadores no se anulan.
\end{corolario}

\addtocounter{definicion}{1}%

\section{Fórmulas de producto a suma y de suma a producto}

\begin{proposicion}[Producto a suma]\label{prop:producto-suma}
Para todos $\alpha,\beta$,
\begin{align*}
\cos\alpha\cos\beta&=\tfrac12\bigl[\cos(\alpha-\beta)+\cos(\alpha+\beta)\bigr],\\
\sin\alpha\sin\beta&=\tfrac12\bigl[\cos(\alpha-\beta)-\cos(\alpha+\beta)\bigr],\\
\sin\alpha\cos\beta&=\tfrac12\bigl[\sin(\alpha+\beta)+\sin(\alpha-\beta)\bigr].
\end{align*}
\end{proposicion}

\begin{corolario}[Suma a producto]\label{cor:suma-producto}
Para todos $x,y$,
\begin{align*}
\sin x+\sin y&=2\sin\tfrac{x+y}{2}\cos\tfrac{x-y}{2}, & \sin x-\sin y&=2\cos\tfrac{x+y}{2}\sin\tfrac{x-y}{2},\\
\cos x+\cos y&=2\cos\tfrac{x+y}{2}\cos\tfrac{x-y}{2}, & \cos x-\cos y&=-2\sin\tfrac{x+y}{2}\sin\tfrac{x-y}{2}.
\end{align*}
\end{corolario}

\addtocounter{definicion}{1}%

\addtocounter{definicion}{1}%

